\subsection{Application: new characterizations of discrete valuation rings}

PROPOSITION 11. Let A be a local Krull domain (in particular an integrally closed local Noetherian domain) and m its maximal ideal. The following conditions are equivalent:

\begin{enumerate}
\def\labelenumi{(\alph{enumi})}
\item
  A is discrete valuation ring;
\item
  m is invertible;
\item
  A: m ≠ A;
\item
  m is divisorial;
\item
  m is the only non-zero prime ideal of A.
\end{enumerate}

As every non-zero ideal of a discrete valuation ring is principal (Chapter VI, § 3, no. 6, Proposition 9), it is invertible and hence (a) implies (b). If m is invertible, its inverse is A: m (Chapter II, § 5, no. 6, Proposition 10) and hence A: m ≠ A; hence (b) implies (c). If A: m ≠ A, then A: (A: m) ≠ A; now m ⊂ A: (A: m); hence m = A: (A: m) since m is maximal, so that m is divisorial (no. 1, Proposition 1 (c)); thus (c) implies (d). The fact that (d) implies (e) follows from Theorem 3 of no. 6, Finally, if m is the only non-zero prime ideal of A, it is of height 1 and hence A \(_{m}\) is a discrete valuation ring (no. 6, Theorem 4); as A is local, A \(_{m}\) = A, which shows that (e) implies (a).

